\section{Introduction}

Charts make quantitative relations visible, but they also make scale a moving target. The same visual evidence can be verbalized as a percentage, a decimal, a raw value, or a value in thousands or millions. A chart question answering model must therefore do more than find the relevant bar, line, or label: it must decide which scale the answer should use.

Modern multimodal large language models (MLLMs) achieve strong aggregate accuracy on chart QA benchmarks such as ChartQA \citep{masry2022chartqa}. Aggregate accuracy can nevertheless hide rare failures that are costly when they occur. We study \emph{scale inconsistency}: a model gives a numerically related answer under the wrong scale, such as predicting \textit{0.61} when the chart question requires \textit{61}. This is not a generic numeric miss. The model has often found the right quantity but has failed at the last reliability step: choosing the scale in which to express it. In financial, scientific, or policy charts, an answer that is off by a factor of 100 while looking plausible is more dangerous than an obviously wrong one.

The unconditional interventions we test do not solve this. Normalizing all percentages, training an adapter that prefers larger numbers, or asking the model to self-check units each helps some cases and hurts others. A legitimate answer such as \textit{0.0593} should not be expanded; an axis in millions should not be treated like a percent display; and a prediction that is already too large needs the opposite intervention from one that is too small. The central question is therefore not whether a model can \emph{produce} a scaled alternative, but whether a system can \emph{select} the correct scale regime: the direction, unit context, and output convention required by the chart and question.

Natural strict scale errors are sparse under any single model and split, so the evidence must be read at two levels. Across the ChartQA validation/test outputs we audit, strict factor-of-100 candidates account for 0.12--1.09\% of examples (Table~\ref{tab:strict_prevalence}). We treat the automatic ratio rule as a \emph{candidate filter}, not as proof of a scale error, and verify test-set candidates with a blinded protocol in which the annotator sees only the chart, question, gold answer, and prediction (\S\ref{sec:strict}). To understand the mechanism behind these sparse failures, we also construct ScaleQA-Natural, a mined and author-adjudicated set of 520 natural scale-conflict candidates, 473 of which are directly image-evaluable (ScaleQA-Natural-473). It is a diagnostic set, not a replacement for official ChartQA evaluation.

Our main empirical finding is that scale correction is directional and regime-dependent. In an answer-tuned Qwen2.5-VL-7B setting, a prefer-large scale-contrastive adapter improves held-out ChartQA test accuracy by +0.72 points (paired bootstrap 95\% CI $[+0.38,+1.09]$) and recovers \verify{14/17} verified strict test-set scale errors; compute-matched continuation controls with ordinary answers, random numeric negatives, or non-directional pairs recover at most \verify{1/17}, so the effect is attributable to directional scale supervision rather than extra training. The aggregate gain is small because the failure is rare; the targeted effect is large. Yet the same adapter is unsafe when applied globally: on ScaleQA-Natural it breaks more examples than it fixes, and the best direction reverses across subtypes. A blind-generation analysis shows the model often already produces the correctly scaled answer (pass@20 \verify{79.3\%} versus \verify{52.4\%} single-sample), while label-free selectors recover only a small part of that headroom. Candidate generation and regime selection both matter, and regime selection remains the larger open gap.

\noindent\textbf{Contributions.}
We contribute (i) a verified strict scale-error recovery probe on ChartQA test and ScaleQA-Natural, an audited diagnostic set of natural scale conflicts;
(ii) evidence, with compute-matched controls, that directional scale supervision repairs a verified error slice but is unsafe as a global policy;
(iii) a decomposition of candidate generation from regime selection, using prompts, verifiers, calculators, likelihood scoring, blind sampling, and label-free selectors; and
(iv) a scoped conclusion: for percentage--decimal, factor-of-100 errors on ChartQA-like data, reliable scale-regime selection, not candidate availability, is the central unresolved problem.
